\section{Open problems}
\label{ch:open-problems}

This chapter collects the questions the project leaves open.
It starts with the paragraph on open problems of the paper's discussion, reproduced as it stood before the paper was shortened: the kernel of the quartic form above level zero, the six-point exchange and the covering of moduli space, the zero of the level-zero kernel at complex momentum, the conclusion of the mean field, the quantum master equation, the single-strand route and the two-string gauge parameter.
A section written for this document then gives the zero of the level-zero kernel with the numbers behind it.
The sections after it are the open problems that the paper does not state, because no sentence of the paper depends on them; each carries its status and points to the chapter on which it bears.
The chapter ends with a comparison that exists only here, between the cubic theory in level truncation and the cubic action for gravity of Cheung and Remmen.

\subsection{The open problems of the paper}
\label{sec:open:paper}

\status{Open}
The block below is the paragraph on open problems of the paper's discussion section, in the notation of the paper.
Its pointers to the sections of the paper refer to blocks reproduced in other chapters: the kernel of the quartic form above level zero is treated in \cref{ch:momentum-space}, the six-point exchange and the single-strand route in \cref{ch:worldsheet-circle}, the mean field in \cref{ch:solutions}, the master equation and the two-string gauge parameter in \cref{ch:gauge-bv}.
The scan of the level-$(1, 1)$ block at finite momentum is the computation the first problem names; no further section is devoted to it.

\subparagraph{Open problems.}

We now collect in physical terms the questions left open along the way, saying for each what is asked, why it matters for the construction and what would settle it.

The first is the kernel of the quartic form at higher level.
The auxiliary field lives on the quotient of the $2$-string states by the kernel of $\vertex{4}$, and the construction on the quotient assumes, from the quintic order on, either that the kernel is empty or that the quintic and sextic vertices vanish on it.
At level zero, for the tachyon pairs at every real pair momentum, the kernel is empty with Euclidean relative momenta, as in \cref{p:sec:worldsheet}, for the vertex family of~\cite{Erbin:2026:SL2CVerticesClosed} at the sections considered there with $\beta > \sqrt{2}$.
The pair kernel is then strictly definite on the even functions of the relative momentum.
At zero momentum the level-$(1, 1)$ block on the tachyon and the ghost dilaton is invertible (\cref{p:sec:solutions}).
Above level zero, the kernel is absent at generic momentum under the Fredholm property and the physical exclusion stated above.
  The question is whether it is absent at every real momentum.
  None of the mechanisms known to force zeros of string vertices (the dilaton theorem with the ghost-number counting, BRST-exact strands, the universality of the couplings at zero momentum), nor the worldsheet twist, produces a null vector of the quartic form at the levels and momenta where its entries are known.
  However, these mechanisms are not exhaustive.
  Nor can a positivity argument exclude a real zero: for the same vertex family, the level-$(1, 1)$ block is indefinite near zero momentum, so isolated real zeros are not excluded.
  The first block to compute is the level-$(1, 1)$ block at nonzero momentum, of which~\cite{Erbin:2026:SL2CVerticesClosed} now provides the first entries.

The question matters beyond consistency.
  At a real zero of the quartic form with a single null vector, the kernel condition of the quintic order reduces to the vanishing of a single pairing of the null vector with the quintic source at that momentum, a property of $\vertex{5}$ that nothing requires.
  If this pairing does not vanish, the composite vertex of \cref{p:sec:hs:quintic} has a pole in that momentum.
The reduced sextic vertex, $\vertex{6}$ minus six fifths of the composite vertex, then carries the opposite pole.
  The pole cancels in the combination $\dvertex{6} + \tfrac{6}{5} \, \vertex{5 \cdot 5} = \vertex{6}$.
  The pole is a propagating state in the $\Sigma_2$ channel, and it makes the reduced vertices singular there.
A real zero above level zero would be accidental (\cref{p:sec:hs:quartic}).
  We expect the pairing not to vanish on it, so that the multiplier on the kernel directions takes over there.
A computation of the momentum-dependent blocks of $\vertex{4}$ and $\vertex{5}$ at the first levels, from an explicit vertex family, would settle whether the quotient suffices or a multiplier on the kernel has to be used.

The second is the six-point exchange and the covering of moduli space.
The composite vertex $\vertex{5 \cdot 5}$, two quintic vertices joined by a $\Sigma_2$ line, has the dimension of $\mathcal M_{0,6}$.
Whether $\tfrac{6}{5} \vertex{5 \cdot 5}$ is the integral over a sub-region of $\mathcal V_{0,6}$, covered once, is open (\cref{p:sec:worldsheet}).
We conjecture that it is, with $\dvertex{6}$ the integral over the complement, so that the two cover the vertex region once.
  One route to a proof reads the composite vertex as glued surfaces, which requires the inverse of the quartic form to act as a gluing (\cref{p:sec:worldsheet}).
  For the explicit section of~\cite{Erbin:2026:SL2CVerticesClosed}, the inverse, summed over levels, does not act as a gluing with a Möbius cut coordinate, within the classes of gluings and of pair states specified there.
  Gluings with non-Möbius cut coordinates remain open.
At fixed level and genus zero, the results of \cref{p:sec:worldsheet} constrain the answer through the poles of the composite vertex in the pair momentum.
  A pole away from the plumbing lattice (the poles of the plumbing fixture) would rule out a compact sub-region with momentum-independent boundary and the standard measure, and the momentum-dependent regions of sector type considered there.
  We have neither established nor excluded such a pole.
  With a pole off the mass shells, on the plumbing lattice or not, or with no pole at all, the composite vertex is not a union of Feynman regions either, modulo the standard factorization theorem.
  The case of poles on the mass shells only is not settled.
  With a pole off the plumbing lattice, the object that remains is a momentum-weighted density on a region, such as a cusp-shaped region pinched at a boundary point, which can carry it.
  Whether such a region reproduces the composite vertex is open.

The question matters because it decides whether the reduced sextic vertex is itself a region of moduli space or an algebraic combination only, that is, what the exchange of $\Sigma_2$ corresponds to on moduli space at six points.
At fixed level, whether the composite vertex has a pole off the plumbing lattice is decided by computing the complex zero locus of the determinant of the quartic form in the pair momentum, with the residue at each zero.
If there are poles, matching their residues against a cusp-shaped region would decide at that level whether such a region reproduces the composite vertex; if there are none, the sub-region question stays open.
This computation needs the momentum-dependent blocks of the previous paragraph.
Even then, the fixed-level composite vertex is not the truncation of the full one, which inverts the quartic form on all $2$-string states at fixed external level.
A zero found at a fixed level is evidence only if it persists as the level grows.

The third is the level-zero zero.
At complex pair momentum, where a pole of the composite vertex off the plumbing lattice would sit, if there is one, the determinant of the discretized level-zero kernel vanishes at a point off the plumbing lattice.
This zero is stable under the discretization (\cref{p:sec:worldsheet}).
Whether the kernel itself has a square-integrable null function there is open.
The null vector of the discretized kernel is not resolved at the precision reached, and deciding it needs the kernel's entries to a precision well beyond it.
If a null function exists, its coupling to the quintic vertex decides at level zero whether the composite vertex has the pole of the previous paragraph.
As long as this is open, the pole hypothesis is neither realized nor excluded at level zero.
Whether a level-zero zero survives in the full composite vertex is the question of the level limit above.

The fourth is the conclusion of the mean-field approximation.
In the quantum theory the exact equation of the auxiliary field at quartic order says that $\Sigma_2$ is, formally, the pair correlator of the string field in its own background (\cref{p:sec:solutions}).
The Gaussian approximation to it gives a gap equation, as in the Hubbard--Stratonovich derivation of the BCS gap equation.
In the level-$(1, 1)$ instance at zero momentum, with the cubic vertex dropped, we solved that gap equation completely.
At the values of the coupling we examined, no root is a minimum of the potential on the real slice.
  Among the real roots, the only local convergent saddles on an admissible contour of the rotated $\Sigma_2$ integral are the one-loop displacements of the classical point.
A mean-field treatment thus gives back the classical structure and no new vacuum.

This computation has several limits.
  The mean field keeps the direct channel, one third of the one-loop self-energy (Hartree and Fock terms together), and drops the diagrams with an internal $\Sigma_2$ line, which no small parameter suppresses.
The obstacles to a BCS-type condensate are arguments, not theorems.
The instance has no loop momentum, no ghost sector in the trace and no cubic vertex.
We expect its conclusion to persist with loop momentum, but we have not established it.
  We have not examined the ghost sector and the cubic vertex.
The finite-dimensional models of \cref{p:app:finite-models} have no momenta either, and the mass shells and the on-shell behavior are outside them.
The gap equation with loop momentum, the ghost sector and the cubic vertex restored would settle whether the conclusion of the mean field persists.
The diagrams the mean field drops remain beyond this gap equation.

The fifth is the quantum master equation, which we do not address in this paper (\cref{p:sec:gauge}).
The construction is classical: the master action of the cubic theory is that of the original theory plus a trivial pair, through order $\gs^4$.
We do not include the loops of the auxiliary field, starting with the $\Sigma_2$ tadpole on the quintic vertex, nor the loop vertices of the original theory, which do not enter the classical construction.
  In this tadpole, the two $\Sigma_2$ legs of the coupling $\vertex{5}(\Sigma_2^2, \Psi)$ are joined by a $\Sigma_2$ propagator.
  Nor do we treat the Batalin--Vilkovisky structure of the multiplier form, whose multipliers need their own antifields.
  Beyond sextic order, we have not constructed the flat representative for the pair of $\Sigma_3$.
The question matters because the auxiliary field is a $2$-string field with its own propagator.
Whether its loops are absorbed in the loop vertices of the original theory, or generate vertices of their own, is what the quantum master equation would decide.
Settling it needs the master action with the loop vertices of the original theory and the auxiliary loops, and its master equation at one loop, with the trivial pair kept.

The sixth is the single-strand route, a different cubic theory suggested by the circle picture.
It has one asymmetric cubic vertex with a special puncture and an auxiliary field with a single string leg.
The quartic vertex would be produced by the exchange of that field after a change of its kinetic term, as in the open-string stub~\cite{Erbin:2024:OpenStringStub}.
We have established two things, the counting and the covering (\cref{p:sec:worldsheet}).
At four points, with the reversed plumbing range, every ghost and moduli count passes with sphere topology once the local discs of the asymmetric vertex extend beyond the unit circle.
The plumbing family covers the interior of $\mathcal V_{0,4}$ exactly once when the stub length varies along a seam symmetric in the three channels, the seam family of \cref{p:sec:worldsheet}, and never for a constant stub length.
This is numerical, at the parameters sampled, and the seam is under-determined.

Off-shell gauge invariance is not established: the local coordinates transported across the seam agree to first order only, and the strictly cubic single-strand theory is expected, though not shown, to fail gauge invariance off-shell at second order in the seam data on non-primary external states.
The question matters because such a theory would be cubic with a single string field and a single auxiliary field.
Its auxiliary channel would be a gluing region, which the inverse of the quartic form is not, in the class of kernels of \cref{p:sec:worldsheet}.
A check of the off-shell identity at second order, and of the embedding condition under which a quartic term on the collar of the seam restores it, would settle it.

The last is the $2$-string gauge parameter.
We have asked whether the cubic action is the gauge-fixed form of a theory with an independent $2$-string parameter $\Lambda_2$ (\cref{p:sec:gauge}).
  In the classes we have examined, $\Lambda_2$ acts on $\Sigma_2$ alone through a field-independent operator, while $\Psi$ keeps its gauge transformation and $\Sigma_2$ its composite law, on the present space of $2$-string states or on an enlargement of it, including the relaxation of level matching with operators built from the zero modes.
  Every such theory is a trivial enlargement: $\Lambda_2$ shifts $\Sigma_2$ along null directions of the enlarged quadratic form and carries no physics.
  At order $\gs^2$, on the present field content and at generic momentum, the cubic action exposes no other $2$-string symmetry, up to the components that the action does not constrain at this order.
  We do not exclude a theory of which the cubic action is a gauge-fixed form and in which $\Lambda_2$ acts on $\Psi$, or through a field-dependent operator, or through the quartic vertex off level matching.
  On the present field content, a $2$-string symmetry could still appear at order $\gs^3$, which we have not examined.



\subsection{The zero of the level-zero kernel in complex momentum}
\label{sec:open:level-zero-zero}

\status{Open}\record{051-kp-complex-s (PH18 at level 0)}
This section gives the computation behind the third problem of the block.
It works under the Euclidean prescription of \cref{sec:mom:level-0}, with the explicit section of the four-point vertex at $\beta = 2$ and its two inputs on the mapping radii and the tile assembly (\cref{ch:sl2c-vertex}).
In the convention of the computation the squared pair momentum is $s = P^2$, all momenta incoming and the signature mostly plus, and the poles of the plumbing fixture sit on the lattice $\alpha' s \in \{4, 0, -4, -8, \dots\}$.
At level zero there is no determinant of the quartic form: the pair kernel $K_P(s)$ is a trace-class operator in the relative momentum (\cref{ch:momentum-space}), and the level-zero form of the pole question is whether $K_P(s)$ has a square-integrable null function at some complex $s$ off the lattice, and whether that function is isotropic for the bilinear pairing.

\paragraph{A scalar model.}
Freezing both relative momenta at zero reduces the kernel to one number, the contact term $F(s, 0, 0)$, an entire function of $s$, positive on the real axis.
It has simple complex zeros off the lattice: the pair $\alpha' s = - 1.7125875239 \pm 5.7403212597 \, \I$ at $\beta = 2$, and at $\beta = \tfrac{3}{2}$ none in the scanned strip $\Re \alpha' s \in [-40, 20]$, $\abs{\Im \alpha' s} \le 10$, the first pair being $\alpha' s = - 1.7583 \pm 10.9120 \, \I$.
A $1 \times 1$ block has no isotropic null vector, so the mechanism of a pole off the lattice is realized in this model.
The model decides nothing at level zero, since one matrix element of $K_P$ is not an eigenvalue branch; it served to place the search contour.

\paragraph{The discretized kernel.}
On the Nyström matrices of the $s$-wave block of $K_P(s)$ at $\beta = 2$, the determinant has one simple zero off the lattice, at
\begin{equation}
  \alpha' s_{**}
    \approx 3.2231131 + 9.4645234 \, \I .
\end{equation}
Five discretizations place it within a disc of diameter $2.8 \times 10^{-7}$: $44$ and $56$ radial nodes at the radial cutoff $\Lambda = 8$, the cutoff $\Lambda = 10$ at the same node density, and three rules for the vertex integral with $60$, $90$ and $120$ points.
The zero is therefore a property of the discretized family, and neither a cutoff artifact nor a quadrature error.
Whether $K_P(s_{**})$ itself has a square-integrable null function is not established.
No separated eigenvalue branch is seen to vanish there, and the null vector of the matrices is not resolved at the precision reached, about $31$ digits: under every regularization it sits at the edge of the radial range, and its coefficients on the eigenvectors of the kernel do not decay over the resolved range.
Two readings remain: a square-integrable null function whose coefficients lie below the precision floor, or none at all.
What would decide between them is the profile of these coefficients computed from matrix entries accurate beyond $10^{-50}$ relative to the kernel, which needs a vertex-integral rule converged at that level (the present rule is converged to about $10^{-13}$, not monotonically) and an analytic derivative in $s$; this computation was not run.
Only if a null function exists does its isotropy become the question.

\paragraph{The rest of the strip.}
On $\Re \alpha' s \in [-40, -20]$, $\abs{\Im \alpha' s} \le 40$, none of the four leading eigenvalue branches of the $s$-wave and the $d$-wave vanishes (for the fourth branch of the $d$-wave, only on $\Re \alpha' s \in [-30, -20]$, $\abs{\Im \alpha' s} \le 30$); this is strong evidence and not a certificate, since no interior test was run there.
On $\Re \alpha' s \in [-20, 20]$, $0 \le \Im \alpha' s \le 40$, a count of the zeros of the leading branches by their winding finds, besides $s_{**}$, four cells with winding one that the refinement could neither attribute to a zero nor exclude; a winding without a zero is the expected outcome there, a zero the alternative.
The search covers neither higher partial waves, nor $\abs{\Im \alpha' s} > 40$, nor $\beta = \tfrac{3}{2}$.
As long as the null function is not decided, the pole of the composite off the lattice is neither realized nor excluded at level zero, and whether a level-zero zero survives in the full composite is the question of the level limit in the block above.

\subsection{Covering of the six-punctured sphere by patches with three punctures}
\label{sec:open:npatch}

\status{Open}\record{081-npatch-n3 (R90); 090-h1-error}
The patch picture of \cref{sec:wscircle:npatch} builds the patch for $n = 3$, the surface of the symmetric quartic vertex with one leg made special, the map of two such patches glued by plumbing to $\mc M_{0, 6}$, and a decomposition of the condition that the patches cover the moduli space once into necessary pieces.
Several of these pieces are untested, and so are the global injectivity and the surjectivity of the map.
The next steps are known: use the center of the section, not its centroid; look for the locus where the map overlaps itself, not for the zeros of its Jacobian; test the inversion in the own channel with the acceptance $\abs{q'} < \abs{q}$; and exclude from the march the trivial root of the own channel at $\abs{q} = 1$.
The paper states the condition with three or more punctures from the atlas picture only and needs nothing more.

\subsection{The inverse of the quartic form off shell}
\label{sec:open:inverse-offshell}

\status{Open}
Off shell the quartic form is a gluing over the vertex region in which the external coordinates move with the plumbing parameter and the antighost insertions are dressed (\cref{ch:worldsheet-circle}); no deformation of the canonical propagator equals it off shell.
Its inverse is known in closed form only through the fixed-halves model, which is exact on shell.
Open are the joint statement on the circle map beyond the pointwise one checked at output level $K \le 2$ (a joint-rank test at finite level and the condition on the side of the kernel), the matrix element of the BRST operator on the seam of the rigid gluing, and the cure of the neck map.
In the fixed-halves model, the inverse summed over levels is singular at real timelike pair momentum in the even sectors with angular number $m \neq 0$, at $\alpha' P^2 \le - 8.71$ for $m = 4$, where an angular Fourier mode of the weight of the region vanishes; no block truncated in level is singular there (\cref{sec:wscircle:note}).
Whether the actual quartic form shares this singularity has not been examined; the assumption of the paper that the propagator of $\Sigma_2$ has a limit at generic momentum covers it.

\subsection{The quartic vertex when the channels cover but are not dual}
\label{sec:open:not-dual}

\status{Open}
If the regions of the three exchange channels tile $\mc M_{0, 4}$ but the local coordinates of two channels disagree on their common edge, BRST invariance still needs a quartic vertex.
It would be a degenerate chain of the space $\mc P_{0, 4}$ of spheres with local coordinates, lying above the boundary lines: a one-parameter family of local coordinates interpolating between the two channels, of zero area in $\mc M_{0, 4}$, so that the vertex vanishes on shell and not off shell.
This argument is ours and has not been checked against the literature; nothing in this document relies on it.
A reading of the theorem of Sonoda and Zwiebach~\cite{Sonoda:1990:CovariantClosedString} confirms its premise, since with the domain of the local coordinates left free every symmetric Hermitian vertex covers and fails on the seams, but it does not address the form of the quartic vertex.

\subsection{The explicit quartic region and the classical constructions}
\label{sec:open:classical-regions}

\status{Open}
Three comparisons between the quartic region of the explicit vertex family (\cref{ch:sl2c-vertex}) and the constructions of the late 1980s remain to be made; none is needed by the paper.
The first is whether the local coordinates of the explicit region derive from a quadratic differential in the sense of~\cite{Sonoda:1990:ClosedStringField}, and whether the interpolation used there can be written for this family.
The second is a numerical comparison of the explicit region with the restricted tetrahedron of~\cite{Saadi:1989:ClosedStringField}, through the Strebel coordinates $(a, b, c)$ of the four-punctured sphere.
The third is the one-loop analogue of the circle picture of \cref{ch:worldsheet-circle}, against the conjecture of~\cite{Sonoda:1990:CovariantClosedString} that no symmetric vertex gives a modular-invariant one-loop tadpole.

\subsection{The hypotheses on the kinetic operator above level zero}
\label{sec:open:above-level-zero}

\status{Open (PH35, PH36)}
The absence of a forced null vector of the kinetic operator of $\Sigma_2$ above level zero, in \cref{sec:mom:fixed-level}, rests on two hypotheses: a Fredholm property of the operator at fixed level, after the level-zero Gaussian is divided out on both sides, and the statement that no pair state is blind to the exchange of single-string states at every mass level.
Each has a named discharge.
The Fredholm property would follow from the symbolic form of the level-$(1, 1)$ block of $\sprod{2}{2}$ and the ellipticity of its weighted symbol matrix.
The decoupling clause of the second hypothesis is a finite computation with the cubic vertex alone, cheapest at level one in the ghost sector.
Two further points are open.
Isolated real zeros of $\sprod{2}{2}$ above level zero are not excluded, since the level-one form is indefinite near $P^2 = 0$; the paper does not need to exclude them, because it states its assumption at generic momentum and the multiplier of \cref{ch:multiplier-hybrid} handles a zero.
And the definiteness at level zero of the kinetic operators of the fields $\Sigma_k$ with $k \ge 3$ needs a sign condition pointwise on the section of the six-point vertex; the evenness of the kernel is proved not to replace it.
The hypotheses are listed in \cref{ch:hypotheses}.

\subsection{Landau analysis and null vectors of composite vertices}
\label{sec:open:landau}

\status{Open}
The question is whether the ideas of Landau analysis, the pinching of integration contours and the endpoint singularities, say something about the null vectors of composite vertices.
A composite vertex that is the integral over a region of moduli space is entire in the momenta, except where the region reaches a degeneration; there, the endpoint of the plumbing integral gives the poles of the plumbing lattice, the string analogue of an endpoint singularity.
A null vector of $\sprod{2}{2}$ at a momentum off the lattice would put into the composite a pole with no geometric origin: either its residue vanishes, which is the kernel condition of \cref{ch:all-orders}, or the composite is not the integral over a region.
The question is therefore tied to the pole of the composite off the lattice in the block above and to the covering conjecture.
The singularity of the inverse summed over levels in the fixed-halves model (\cref{sec:open:inverse-offshell}) lies off the lattice, consistent with the absence of an inverse region in that model.
This sketch has not been reviewed.

\subsection{Routes toward the six-point covering}
\label{sec:open:six-point}

\status{Open}\record{105-exchange-real-momenta (R99, R100); 110-p22-inverse-sewing (R101)}
The block above states the covering of the six-point vertex region by the exchange and the dressed sextic vertex as a conjecture, with the normalization $\dvertex{6} + \tfrac{6}{5} \vertex{5 \cdot 5} = \vertex{6}$.
What is established is listed in \cref{ch:worldsheet-circle} and \cref{sec:ws1:sewing}: the tests at real momenta are necessary conditions, which can refute the conjecture but not prove it; the inverse of the quartic form off shell has no sewing form within kernels that act trivially on the multiplicities of the matter $\mathfrak{sl}(2)$ (the pair of conformal Killing vectors, ghost dressings, Möbius cut coordinates); and the chain property fails at first order on an open subset of each face of $\mc V_{0, 5}$ where a cubic vertex meets a quartic one through a propagator.
The routes left open are a sewing form with cut coordinates that are not Möbius and depend on the plumbing parameter, for which a necessary test at second order is stated and would not decide even if it passed; the shape in which the legs are dressed, for which no closing step is known; the reading as a chain without factorization; and a condition on the kernel of the circle map that can be computed from the gluing recipe of \cref{ch:worldsheet-circle} alone.
None of them decides without a quintic vertex, and the construction of a quintic vertex, from an extension of the explicit family to five points or from the Strebel family of~\cite{Moeller:2007:ClosedBosonicString}, is the step this line of work has not taken.

\subsection{Routes for the two-string gauge parameter}
\label{sec:open:lambda2}

\status{Open}
The last problem of the block is treated in \cref{sec:gauge-bv:unfixing}, whose final section lists the classes outside the exclusions (\cref{sec:gauge-bv:unfixing-open}).
Three routes are not there.
The first is the dressing at order $\gs^3$ in a finite model: extending the solve of order $\gs^2$ to the quintic term $\vertex{5}(\Sigma_2^2, \Psi) / 5!$ and closing the algebra, which needs a quintic vertex consistent with the \Loo{} relations of the model at that order; it is the only route on the present field content, and a computation in a model, not on worldsheet data.
The second is on the \IBL{} side (\cref{ch11:sec:mc}): a modification of the products $\sprod{m}{n}$, compatible with the \IBL{} relations, that removes the mismatches of the Maurer--Cartan form, and a re-identification of $\gs^2$ with the loop parameter of the \IBL{} structure, or a reading of $\Sigma_2$ as a one-loop object, under which the bubble term produced by the coupling $\vertex{3}(\Sigma_2, \Psi)$ that the Maurer--Cartan form requires sits at loop order.
For that bubble it is open whether a linear redefinition independent of $\gs$ absorbs its term of order $\gs^0$ for the closed string, which a finite model excludes, and whether the bubble is finite and nonzero, which a level truncation or a gauge test would decide.
The third is a worldsheet definition of the two-string form $q$ from one factorization channel of the four-punctured sphere: such a form need not be totally symmetric in its four slots, so it would carry a component that vanishes on every square $\Psi^2$, invisible in the eliminated action at order $\gs^2$: a form $q$ that carries data beyond $\vertex{4}$, the case left open in \cref{sec:gauge-bv:projector}.

\subsection{Further directions}
\label{sec:open:directions}

\status{Open}
Three directions were proposed and not pursued.
Multilocal fields in momentum space, $\sigma_n(k_1, \dots, k_n)$, could be treated with the Wigner--Weyl transform, as in condensed-matter systems with inhomogeneous backgrounds; the circle map of \cref{ch:worldsheet-circle} already encodes the relative momenta in the level content.
The worldsheet description of the fields could use \Loo{} algebroids instead of the coloured operad of \cref{ch:worldsheet-dictionary}.
And a cubic closed string field theory could perhaps be built from the open--closed double copy, by analogy with the cubic action for gravity built on the symmetry of the double copy.

\subsection{Comparison with the cubic action for gravity of Cheung and Remmen}
\label{sec:open:cheung-remmen}

\status{no ledger entry}\record{099-cheung-remmen-comparison}
Cheung and Remmen wrote the Einstein--Hilbert action, to all orders around flat space and in $D \ge 3$, as a cubic action of the metric and an auxiliary field~\cite{Cheung:2017:HiddenSimplicityGravity}.
This section asks whether the cubic closed string field theory, restricted to a level truncation that keeps the massless sector, can be matched to that action.
The answer is no: a field-redefinition invariant of the massless sector differs between the two at every finite level.
The two share a mechanism, but only at quartic order and only for the Einstein--Hilbert action itself, not for the truncated string field theory; at the level of an analogy, the two constructions are of the same kind, as is the strictification of Yang--Mills theory of \cref{ch:qft-toys}.
The results of Cheung and Remmen and of Yang and Zwiebach used below are cited, not re-derived; the algebra of the second part is checked in a finite-dimensional model by a script; the comparison as a whole has not been reviewed independently.

\paragraph{Setting and notions of match.}
On the string side, $\Psi$ is restricted to the states of level at most $L$, with $L \ge (1, 1)$ so that the massless sector (graviton, $B$-field, ghost dilaton, and the auxiliary vectors outside Siegel gauge) is kept, and $\Sigma_2$ lives on the symmetric $2$-string states of the truncated space modulo the kernel of $\vertex{4}$; the vertices are those of the full theory restricted to the kept states, with no massive field integrated out.
The quartic step of the construction is
\begin{equation}
  \label{eq:open:cr:hs4}
  \widetilde S_4
    = - \frac{\gs^2}{4!} \, \vertex{4}(\Sigma_2^2)
    + \frac{2 \gs^2}{4!} \, \vertex{4}(\Sigma_2, \Psi^2).
\end{equation}
On the gravity side, with the densitized inverse metric $\sigma^{ab} = \sqrt{- g} \, g^{ab} = \eta^{ab} - h^{ab}$ in mostly-plus signature, the action of Cheung and Remmen is
\begin{equation}
  \label{eq:open:cr:action}
  \mc L_{\rm EH}
    = - \left(
      A^a{}_{bc} A^b{}_{ad}
      - \frac{1}{D - 1} \, A^a{}_{ac} A^b{}_{bd}
    \right) \sigma^{cd}
    + A^a{}_{bc} \, \pd_a \sigma^{bc},
\end{equation}
in their normalization, with an overall $1 / 16 \pi G$ omitted.
A field redefinition here is perturbative, $\Phi' = L \Phi + O(\Phi^2)$ with $L$ invertible and possibly non-local, and auxiliary fields are added or removed only by their algebraic elimination.
Three notions of match are distinguished, from the strongest to the weakest:
\begin{enumerate}
  \setlength{\itemsep}{0.5ex}

  \item a match of fields and actions: the fields $(h_{ab}, A^a{}_{bc})$ and the action~\eqref{eq:open:cr:action} arise from the truncated cubic theory after a field redefinition;
  \item a match of mechanism: in both, an algebraic auxiliary field coupled linearly to a composite quadratic in the fields, with a field-independent invertible kernel, replaces a quartic contact term by its exchange, the same term in both;
  \item an analogy: both rewrite a non-polynomial gauge theory as a cubic one with algebraically eliminable auxiliary fields.
\end{enumerate}

\paragraph{No match of fields and actions.}
The obstruction is a flat direction.
Let $a$ be the coefficient of the zero-momentum marginal state $c \bar c \, \pd X \bar \pd X$ along a compact direction $X$, a constant deformation of $G_{XX}$.
For every constant $a$ the metric $g = \eta + a \, \delta_X \delta_X$ is flat, so the Einstein--Hilbert action, a two-derivative functional, vanishes on the whole family and every member solves the equations of motion: the zero-momentum potential of $a$ vanishes to all orders, and so it does for the action~\eqref{eq:open:cr:action}, whose equations of motion are equivalent to Einstein's.
In closed string field theory truncated at level $\ell$ (the level of~\cite{Yang:2005:TestingClosedString}, $\ell = L_0 + \bar L_0 + 2$), the quartic term of the effective potential of $a$ is
\begin{equation}
  \label{eq:open:cr:potential}
  \kappa^2 \, \mathbb V_{(\ell)}(a)
    = \bigl( \mc C(\ell) + I_4 \bigr) \, a^4
    + O(a^6),
  \qquad
  \mc C(\ell)
    = \sum_{\ell' \le \ell} c(\ell'),
\end{equation}
with $c(\ell')$ the exchange of the massive fields of level $\ell'$ through the cubic vertex and $I_4$ the elementary quartic vertex~\cite{Yang:2005:TestingClosedString}.
Their values are $I_4 = - 0.25598$ and $c(0) = 0.71191$ from the tachyon, and $\mc C(\infty) = 0.25585$ by extrapolation, the cancellation $\mc C(\infty) + I_4 = 0$ holding to five digits with the value of $I_4$ computed by Moeller.
The coefficient $\mc C(\ell) + I_4$ decreases in magnitude with the level, is nonzero at every level computed, up to $\ell = 300$, and is fitted by $d_0 + d_2 / \ell^2 + d_3 / \ell^3$, which vanishes only in the limit.
A truncation that keeps only the massless sector gives $I_4 \neq 0$, and one that adds the tachyon gives $c(0) + I_4 \simeq 0.456$; the same paper repeats the test for the constant metric and $B$-field in two directions, with the same outcome.
The normalization of~\cite{Yang:2005:TestingClosedString} ($\alpha' = 2$ and a different normalization of the sphere correlator) differs from the one of this document, which changes these numbers by positive factors and not whether they vanish.
The data do not exclude an accidental zero at one finite level, which would remove the obstruction at that level and at quartic order only.
The dilaton direction of the tachyon--dilaton truncation shows the same phenomenon in this project: it is lifted at quartic order by exactly the dropped exchange of level $(2, 2)$, and at quintic order by a dropped exchange as well, under the hypotheses of that computation (\cref{sec:sol:dilaton-d5}).

The Hubbard--Stratonovich step does not change this: at quartic order the elimination of $\Sigma_2$ is exact and gives back the truncated quartic action (\cref{ch:order-by-order}), so the truncated cubic theory and the truncated non-polynomial theory have the same quartic potential at zero momentum.
The obstruction is invariant under field redefinitions.
Let $F$ be a redefinition from the fields of Cheung and Remmen to the truncated fields, with the auxiliary fields of both sides eliminated.
It maps the family $\Phi(a)$ of flat metrics to a family $F(\Phi(a))$ of zero-momentum solutions of the truncated theory, whose tangent $L \Phi'(0)$ at $a = 0$ is an on-shell zero-momentum state of the truncated free theory: in the massless sector, the marginal state up to gauge.
The quartic term of the truncated potential along that state is~\eqref{eq:open:cr:potential} at the truncation level, nonzero, so the truncated equations of motion have no solution tangent to it at order $a^3$, and no such $F$ exists.
The argument needs only that $F$ is perturbative and invertible at linear order, so non-local redefinitions are included, and it does not use the momentum-dependent four-graviton amplitude.
The reason behind it is that the quartic coupling of the Einstein--Hilbert term is the effective coupling of the massless fields after every massive field, the tachyon included, has been integrated out (shown for the superstring in~\cite{Sen:2016:WilsonianEffectiveAction}), followed by the low-energy limit and a field redefinition, while a finite truncation keeps the elementary vertex and a finite part of the exchanges.

A second obstruction holds under a hypothesis: that the block of the quartic kernel on massless pairs, $K_P(q', q)$ at fixed total momentum, has infinite rank as a kernel in the relative momentum.
The motivation is that the Koba--Nielsen factor of $\vertex{4}$ carries $\exp(\alpha' q \cdot q' \, \ell(\xi))$, with $\ell$ a logarithm of cross-ratios that varies over the two-dimensional vertex region, so that the coefficients of the expansion in $q \cdot q'$ are moments of $\ell$ over the region, and the kernel is smooth and not a delta function times a polynomial (\cref{ch:momentum-space}); the Taylor coefficients of the massless block in $q \cdot q'$ would discharge or falsify it.
Under this hypothesis $\Sigma_2$ restricted to massless pairs has infinitely many components at each total momentum, while the auxiliary field of Cheung and Remmen has $D^2 (D + 1) / 2$, which is $40$ in $D = 4$, and a field redefinition, which must be a bijection between the auxiliary spaces at each total momentum, cannot exist.
Level truncation truncates the oscillator content and not the dependence on $q$, so this persists at every level.
The verdict does not rest on this second obstruction.

\paragraph{A match of mechanism at quartic order.}
Write the action~\eqref{eq:open:cr:action} with $\sigma^{cd} = \eta^{cd} - h^{cd}$ as
\begin{equation}
  \label{eq:open:cr:split}
  \mc L
    = - A \cdot M(h) \cdot A
    - A \cdot b(h),
  \qquad
  M(h)
    = M_0 - M_1(h),
\end{equation}
where $A \cdot M \cdot A$ is the quadratic form in brackets in~\eqref{eq:open:cr:action}, split into its constant part $M_0$ and its part $M_1(h)$ linear in $h$, both algebraic, and $A \cdot b(h) = A^a{}_{bc} \, \pd_a h^{bc}$ comes from $A^a{}_{bc} \pd_a \sigma^{bc} = - A^a{}_{bc} \pd_a h^{bc}$.
Products and inverses of $M_0$ and $M_1$ act pointwise on the $D^2 (D + 1) / 2$ components of $A$; $b$ carries the only derivative.
Eliminating $A$ is a Gaussian with a field-dependent kernel,
\begin{equation}
  \label{eq:open:cr:eliminate}
  A_*
    = - \frac{1}{2} \, M(h)^{-1} b(h),
  \qquad
  \mc L_{\rm eff}
    = \frac{1}{4} \, b \, M(h)^{-1} b
    = \sum_{n \ge 2} \frac{1}{4} \, b \bigl( M_0^{-1} M_1(h) \bigr)^{n - 2} M_0^{-1} \, b,
\end{equation}
the $n$-th term being of degree $n$ in $h$: the whole tower of Einstein--Hilbert vertices in this basis is the geometric series of $M(h)^{-1}$, a property of the Einstein--Hilbert action in the basis of the densitized inverse metric and not of a generic non-polynomial action.
On shell $A$ is linear in $h$ at leading order, a combination of Christoffel symbols, while $\Sigma_2 = \Psi^2$ is quadratic.
This mismatch is removed by the local, linear and invertible shift by the linearized solution,
\begin{equation}
  \label{eq:open:cr:shift}
  A
    = A' - \frac{1}{2} \, M_0^{-1} b(h).
\end{equation}
The term $- 2 A' M_0 (- \tfrac{1}{2} M_0^{-1} b)$ cancels $- A' \cdot b$, the cross term from $M_1$ gives a coupling of $A'$ to a bilinear composite, and the action becomes
\begin{equation}
  \label{eq:open:cr:shifted}
  \mc L
    = \frac{1}{4} \, b M_0^{-1} b
    + \frac{1}{4} \, b M_0^{-1} M_1(h) M_0^{-1} b
    - A' \cdot \bigl( M_0 - M_1(h) \bigr) \cdot A'
    - A' \cdot J(h, h),
  \qquad
  J(h, h)
    := M_1(h) \, M_0^{-1} b(h).
\end{equation}
The terms in $h$ alone stop at cubic order, $A'$ has the constant kernel $M_0$ and couples linearly to the local composite $J(h, h)$ and through the vertex $A' M_1(h) A'$, and on shell $A'$ is quadratic in $h$.
The degree-four term of~\eqref{eq:open:cr:eliminate}, with $M_0^{-1}$ and $M_1$ symmetric, is the quartic Einstein--Hilbert vertex in this basis:
\begin{equation}
  \label{eq:open:cr:quartic}
  \mc L^{(4)}_{\rm eff}
    = \frac{1}{4} \, b M_0^{-1} M_1 M_0^{-1} M_1 M_0^{-1} b
    = \frac{1}{4} \, J(h, h) \, M_0^{-1} J(h, h).
\end{equation}
The kernel $M_0$ is non-degenerate and indefinite: its inertia is $(10, 8)$, $(22, 18)$ and $(40, 35)$ in $D = 3$, $4$ and $5$, with determinants $\tfrac{1}{4}$, $\tfrac{256}{81}$ and $- 243$.

Now apply the quartic step~\eqref{eq:open:cr:hs4}, couplings stripped, to a quartic form that factorizes as $\vertex{4}(\Sigma, \Sigma') = \tfrac{1}{4} \, J(\Sigma) M_0^{-1} J(\Sigma')$, with $J$ extended linearly to symmetric $2$-tensors.
Substituting $A' = - \tfrac{1}{2} M_0^{-1} J(\Sigma_2)$ gives $- A' M_0 A' = - \tfrac{1}{4} J(\Sigma_2) M_0^{-1} J(\Sigma_2)$ and $- A' \cdot J(h, h) = \tfrac{1}{2} J(\Sigma_2) M_0^{-1} J(h, h)$, that is,
\begin{equation}
  \label{eq:open:cr:map}
  - \vertex{4}(\Sigma_2, \Sigma_2)
  + 2 \, \vertex{4}(\Sigma_2, h^2)
    = - A' M_0 A'
    - A' \cdot J(h, h)
  \quad
  \text{at}
  \quad
  A'
    = - \frac{1}{2} \, M_0^{-1} J(\Sigma_2),
\end{equation}
which is the quartic-order part of~\eqref{eq:open:cr:shifted}.
When $M_0^{-1}$ is non-degenerate on the image of $J$, the kernel of this $\vertex{4}$ is the kernel of $J$, so the quotient of the symmetric tensors by $\ker \vertex{4}$ is the image of $J$, a space of local fields.
So, after the shift~\eqref{eq:open:cr:shift}, the action of Cheung and Remmen is a Hubbard--Stratonovich rewriting of the quartic Einstein--Hilbert vertex~\eqref{eq:open:cr:quartic}, with a local auxiliary field, a constant non-degenerate indefinite kernel and a linear coupling to a bilinear composite, and the quartic step of this document, applied to that vertex, gives the same field and the same action at quartic order.

This is not a match in level truncation, because the two auxiliary fields replace different quartic terms.
In a massless truncation of the string field theory, $\Sigma_2$ restricted to massless pairs replaces the elementary vertex $\vertex{4}$ restricted to massless states, whose kernel in the relative momentum is smooth and non-polynomial.
The coupling that~\eqref{eq:open:cr:quartic} represents is the effective one: $\vertex{4}$ on massless states plus the exchanges of all the massive fields, which the cubic theory keeps as exchange diagrams of $\Psi$ through $\vertex{3}$ and does not absorb into $\Sigma_2$.
The identification~\eqref{eq:open:cr:map} holds only if the construction is applied to the low-energy effective theory of the massless fields, and only there does $\Sigma_2$ become local, since the quotient by $\ker \vertex{4}$ collapses the bilocal field onto the image of $J$ when $\vertex{4}$ is local.
Beyond quartic order the question is not decided.
Cheung and Remmen use one auxiliary field whose kernel $M_0 - M_1(h)$ depends on $h$, and the quintic step of the construction also makes the kernel of $\Sigma_2$ depend on $\Psi$ (\cref{ch:order-by-order}), the same device; whether the tower of the construction, applied to the Einstein--Hilbert $\vertex{5}$, $\vertex{6}$, \ldots{} in this basis, reproduces $A' M_1(h) A'$ and stops at $\Sigma_2$ is not checked, since the dressing is built from the fully symmetrized $\vertex{5}$ while $M_1(h)$ is one particular split of it.

The script checks the shift~\eqref{eq:open:cr:shifted} as an exact polynomial identity in a model with four components of $A$, three of $h$ and random integer $M_0$, $M_1$, $b$; the Taylor coefficients of $\tfrac{1}{4} b M(t h)^{-1} b$ for $n = 2, \dots, 7$ against~\eqref{eq:open:cr:eliminate} and the quartic one against~\eqref{eq:open:cr:quartic}; the inertia and the determinants of $M_0$, built from~\eqref{eq:open:cr:action} at $\sigma = \eta$, in $D = 3, 4, 5$; and~\eqref{eq:open:cr:map} as a polynomial identity, with $\rank \vertex{4} = \rank J = 5$ on the $10$-dimensional symmetric square of $\R^4$.
Each check has a broken counterpart that fails.
The model checks the algebra of the shift and of the map, not the tensor structure of $M_1$ and $b$; only $M_0$ is the actual one.

\paragraph{An analogy.}
Both constructions make a non-polynomial gauge theory cubic by adding fields eliminated algebraically, with an invertible and indefinite kernel: inertia $(22, 18)$ for $M_0$ in $D = 4$, and an indefinite quadratic form of the auxiliary field already in the tachyon--dilaton truncation of \cref{ch:solutions}.
In both the elimination is exact at quartic order; the field of Cheung and Remmen is algebraic and carries no physical state, and $\Sigma_2$ carries none through the order to which the gauge structure of \cref{ch:gauge-bv} is established.
The differences are as important.
The auxiliary field of Cheung and Remmen is local and also replaces every higher vertex through its $h$-dependent kernel, a property of the Einstein--Hilbert action in this basis; $\Sigma_2$ is bilocal, replaces the elementary quartic vertex only, and the higher vertices need the tower $\Sigma_k$.
In the words of Cheung and Remmen, their construction is a field redefinition away from the first-order Palatini formalism, and the Hubbard--Stratonovich reading needs the shift~\eqref{eq:open:cr:shift}.
The analogy is of the same kind as the one with the strictification of Yang--Mills theory~\cite{Borsten:2021:DoubleCopyHomotopy}, which is an illustration of the mechanism and not a match in level truncation either.

Nützi and Reiterer write the field equations of general relativity around Minkowski space as the Maurer--Cartan equation of a differential graded Lie algebra built from a frame and a Lorentz connection~\cite{Nutzi:2018:ScatteringAmplitudesYM}.
Their variables are first order, there is no auxiliary field replacing a composite, and there is no action: a cyclic pairing of the kind a classical BV action needs would pair degree $1$ with degree $2$ and degree $0$ with degree $3$, while their graded dimensions are $(10, 40, 50, 24, 4)$ in degrees $0$ to $4$.
Their construction strictifies general relativity at the level of the equations, and it bears on the comparison only as a second instance of the analogy.
